\documentclass{amsart}
\usepackage{amsmath,amssymb,amsthm,hyperref}
\usepackage{algorithm}
\usepackage{algpseudocode}
\newtheorem{theorem}{Theorem}
\newtheorem{lemma}{Lemma}
\newtheorem{proposition}{Proposition}
\theoremstyle{definition}
\newtheorem{definition}{Definition}
\newtheorem{remark}{Remark}

\begin{document}

\title{Worst-case complexity of collection in polycyclic groups}
\maketitle

\section{Statement of the result}

We work with a consistent polycyclic presentation of a polycyclic group $G$ of
Hirsch length $n$, given by generators $g_1,\dots,g_n$, relative orders
$r_i\in\{2,3,\dots\}\cup\{\infty\}$, power relations $g_i^{r_i}=w_{i,i}$ (for finite
$r_i$) and conjugation relations $g_j^{g_i}=w_{i,j}$, $g_j^{g_i^{-1}}=w'_{i,j}$ for
$1\le i<j\le n$, each right-hand side being a \emph{normal word}
$g_{i+1}^{e_{i+1}}\cdots g_n^{e_n}$ with $0\le e_k<r_k$ whenever $r_k<\infty$. We use
the convention that the normal form lists generators in \emph{increasing} index from
left to right,
\[
g_1^{a_1}g_2^{a_2}\cdots g_n^{a_n},\qquad 0\le a_i<r_i\ (r_i<\infty),\ a_i\in\mathbb Z\ (r_i=\infty).
\]

Throughout, the cost model is the bit (Turing-machine) model. Let
\[
b \;=\; \text{the maximum bit-length of any integer exponent}
\]
occurring in the relator words $w_{i,i},w_{i,j},w'_{i,j}$ \emph{or} in the input word,
and let $M=2^{b}$, so that every individual exponent in the relators and in the input has
absolute value $<M$. The input is a word
\[
W \;=\; g_{c_1}^{\,\varepsilon_1}\cdots g_{c_\ell}^{\,\varepsilon_\ell},
\qquad c_t\in\{1,\dots,n\},\ \varepsilon_t\in\mathbb Z\setminus\{0\},
\]
of \emph{syllable length} $\ell$ (the number of generator powers written down); the
bit-lengths of the input exponents $\varepsilon_t$ are absorbed into $b$.

\begin{theorem}\label{thm:main}
Let $G$ be given by a consistent polycyclic presentation of Hirsch length $n$ as above,
with relator/input exponent bit-bound $b$, and let $W$ be an input word of length $\ell$.
Then the collection-from-the-left algorithm \textup{(}Algorithm~\textup{\ref{alg:collect})}
computes the normal form of $W$, and:
\begin{enumerate}
\item \textup{(Exponent size.)} Every integer exponent that appears at any moment during the
computation has absolute value at most
\begin{equation}\label{eq:expbound}
B \;:=\; \ell\,M\,(1+nM)^{\,n-1},
\end{equation}
i.e.\ bit-length $\;\log_2 B=O\!\bigl(n\,(b+\log n)+\log\ell\bigr)=O\!\bigl(n\,(b+\log(n\ell))\bigr).$
\item \textup{(Termination and upper bound.)} The algorithm halts, and its running time is
bounded by
\begin{equation}\label{eq:upper}
\ell\cdot 2^{\,O\!\left(n^2(b+\log(n\ell))\right)}\;=\;(nM\ell)^{\,O(n^2)},
\end{equation}
which is polynomial in $\ell$ and in $M=2^{b}$ for each fixed Hirsch length $n$, and
exponential in $n$.
\item \textup{(Lower bound / tightness in $n$.)} The exponential dependence on $n$ in the
\emph{number of elementary collection operations} is unavoidable. There is a family of
consistent polycyclic presentations \textup{(}all relator exponents in $\{-1,0,1\}$, so
$b=1$\textup{)} of Hirsch length $n+1$, and inputs of length $\ell=2n+1=\mathrm{poly}(n)$, on
which Algorithm~\textup{\ref{alg:collect}} performs at least $2^{n}-n$ elementary operations,
i.e.\ a number of operations exponential in the Hirsch length. \textup{(}We emphasise that this
is a lower bound on the \emph{work done by collection-from-the-left}, not on the output size:
the normal form on these instances has $O(n)$ syllables, each exponent a binomial coefficient
of magnitude $\le 2^{n}$, hence of bit-length $\le n$, so the output is of polynomial bit-size
$O(n^2)$. The exponential blow-up is therefore intrinsic to the collection \emph{algorithm}'s
operation count, which expands relator copies one at a time, and not to the size of the
answer.\textup{)}
\end{enumerate}
\end{theorem}

\noindent
Thus the precise answer to ``determine the worst-case complexity'' is
\[
\boxed{\;(nM\ell)^{\,O(n^2)}\;=\;\ell\cdot 2^{\,O(n^2(b+\log(n\ell)))}\;,\qquad M=2^{b},\;}
\]
\emph{polynomial in the word length $\ell$ and in the exponent magnitude $M=2^{b}$ for each
fixed Hirsch length, and exponential in the Hirsch length $n$; the exponential blow-up in $n$
is genuine.}

\medskip
We prove the theorem as follows: \S\ref{sec:alg} specifies the algorithm and its correctness;
\S\ref{sec:exp} proves the exponent bound \eqref{eq:expbound} (this needs no prior knowledge
of termination); \S\ref{sec:steps} bounds the number of elementary operations, which both
proves termination and yields \eqref{eq:upper}; \S\ref{sec:lower} gives the lower-bound
construction.

\section{The algorithm and its correctness}\label{sec:alg}

We represent a group word as a finite list of \emph{syllables}, each syllable a pair $(c,e)$
standing for $g_c^{\,e}$ with $c\in\{1,\dots,n\}$ and $e\in\mathbb Z\setminus\{0\}$. The
\emph{level} of $(c,e)$ is its generator index $c$. A list $(c_1,e_1),\dots,(c_m,e_m)$ is
\emph{collected} (in normal form) if $c_1<c_2<\cdots<c_m$ and, for each $t$, $0\le
e_t<r_{c_t}$ whenever $r_{c_t}<\infty$.

The algorithm keeps a \emph{collected prefix} $P$ (always in normal form, with maximal level
$\mathrm{top}(P)$, taken $0$ when $P$ is empty) and an \emph{uncollected tail} $U$ (a list),
and repeatedly processes the first syllable $(c,e)$ of $U$.

\begin{algorithm}[h]
\caption{Collection from the left}\label{alg:collect}
\begin{algorithmic}[1]
\Require A consistent polycyclic presentation $(g_i,r_i,w_{i,i},w_{i,j},w'_{i,j})$ and a
word $W$ as a syllable list.
\Ensure The normal form of $W$.
\State $P\gets$ empty list;\quad $U\gets W$.
\While{$U$ is nonempty}
  \State let $(c,e)$ be the first syllable of $U$; remove it from $U$.
  \If{$e=0$} \State \textbf{continue} \EndIf
  \If{$c\le\mathrm{top}(P)$}\label{ln:swap}
     \State let $(k,f)$ be the last syllable of $P$, so $k=\mathrm{top}(P)\ge c$; remove it
     from $P$.
     \If{$c=k$} \State prepend $(c,f+e)$ to $U$ \Comment{merge}
     \Else \Comment{$c<k$: move one $g_c$ leftwards}
        \State let $\sigma\in\{+1,-1\}$ be the sign of $e$;
        \State prepend to $U$, in order: $(c,\sigma)$, then the word $\mathrm{Conj}(g_k^{f},g_c^{\sigma})$,
        then $(c,e-\sigma)$.\label{ln:swapprepend}
     \EndIf
  \ElsIf{$r_c=\infty$ \textbf{or} $0\le e<r_c$}\label{ln:place}
     \State append $(c,e)$ to the end of $P$ \Comment{$(c,e)$ is now correctly placed}
  \Else \Comment{$c>\mathrm{top}(P)$ but $e$ out of range; reduce the power}\label{ln:power}
     \State write $e=q\,r_c+e_0$ with $0\le e_0<r_c$ (so $q\ne0$);
     \State prepend to $U$, in order: $(c,e_0)$, then $q$ copies of $w_{c,c}$ if $q>0$, or
     $|q|$ copies of $w_{c,c}^{-1}$ if $q<0$.\label{ln:powerprepend}
  \EndIf
\EndWhile
\State \Return $P$.
\end{algorithmic}
\end{algorithm}

Here $\mathrm{Conj}(g_k^{f},g_c^{\sigma})$ denotes the word obtained from $g_c^{-\sigma}\,g_k^{f}\,g_c^{\sigma}$
by substituting, in each of the $|f|$ letters $g_k^{\pm1}$, the relator
$g_k^{\,g_c^{\sigma}}$; explicitly it is the concatenation of $f$ copies of $w_{c,k}$ (if
$\sigma=+1$, $f>0$), of $|f|$ copies of $w_{c,k}^{-1}$ (if $\sigma=+1$, $f<0$), or with
$w_{c,k}$ replaced by $w'_{c,k}$ when $\sigma=-1$. By the defining property of a polycyclic
presentation, every $w_{c,k},w'_{c,k}$ is a normal word in $g_{c+1},\dots,g_n$; likewise
$w_{c,c}$ is a normal word in $g_{c+1},\dots,g_n$. The formal inverse $w^{-1}$ of a word $w$
is the reversed word with each syllable exponent negated. We record the structural fact used
throughout.

\begin{lemma}[Level monotonicity]\label{lem:level}
Consider any single iteration of Algorithm~\ref{alg:collect} acting on a level-$c$ syllable.
Then:
\begin{enumerate}
\item All syllables in the words $\mathrm{Conj}(g_k^{f},g_c^{\sigma})$ \textup{(}swap with
$c<k$\textup{)} and in the copies of $w_{c,c}^{\pm1}$ \textup{(}power reduction\textup{)} have
level $\ge c+1$.
\item Apart from those words, the only syllables created have level $c$ \textup{(}namely the
leading $(c,\sigma)$ and residual $(c,e-\sigma)$ in a swap, or the residual $(c,e_0)$ in a
power reduction, or the merged $(c,f+e)$\textup{)}.
\end{enumerate}
In particular, no iteration acting on a level-$\ge i$ syllable ever creates a syllable of
level $<i$.
\end{lemma}

\begin{proof}
Immediate from lines~\ref{ln:swapprepend} and \ref{ln:powerprepend} and the definition of
$\mathrm{Conj}$, since every $w_{c,k},w'_{c,k},w_{c,c}$ is a normal word in
$g_{c+1},\dots,g_n$. The last sentence follows because an iteration on a level-$i$ syllable is
a placement (creates nothing), a power reduction at level $i$ (creates levels $\ge i+1$ and a
level-$i$ residual), a merge at level $i$ ($k=i$, creates one level-$i$ syllable), or a swap
with $c=i<k$ (creates levels $\ge i+1$ and two level-$i$ syllables); none produces level $<i$.
\end{proof}

\paragraph{Correctness.}
Each iteration rewrites $P\cdot U$ to an equal element of $G$. Placement and the various
prepends merely re-bracket a word. The merge uses $g_k^{f}g_c^{e}=g_c^{f+e}$ when $c=k$. The
swap uses the identity $g_k^{f}g_c^{\sigma}=g_c^{\sigma}\,\mathrm{Conj}(g_k^{f},g_c^{\sigma})$,
which is exactly the conjugation relation $g_k^{\,g_c^{\sigma}}=w_{c,k}$ (resp.\ $w'_{c,k}$)
applied $|f|$ times; consequently $g_k^{f}g_c^{e}=(c,\sigma)\cdot
\mathrm{Conj}(g_k^{f},g_c^{\sigma})\cdot(c,e-\sigma)$ as group elements. The power reduction
uses $g_c^{r_c}=w_{c,c}$, hence $g_c^{e}=g_c^{e_0}(g_c^{r_c})^{q}=g_c^{e_0}w_{c,c}^{q}$. Thus
the element of $G$ represented by $P\cdot U$ is invariant. When the loop ends, $U$ is empty and
$P$ is collected, so $P$ is \emph{a} normal word equal to the input element. Consistency of the
presentation guarantees the normal word representing a given element is unique; hence $P$ is
\emph{the} normal form of $W$. It remains to bound exponents (\S\ref{sec:exp}) and prove the
loop terminates with the stated number of iterations (\S\ref{sec:steps}).

\section{Exponent size bound}\label{sec:exp}

For a configuration $(P,U)$ and a level $c\in\{1,\dots,n\}$ define the \emph{level-$c$ mass}
\[
N_c\;=\;\sum_{(c,e)\in P\cup U}|e|.
\]
The bound below is proved by causal induction on levels and does not assume termination.

\begin{proposition}\label{prop:exp}
Define $A_1=\ell M$ and $A_c=\ell M+nM\sum_{j<c}A_j$ for $c\ge2$. Then at every moment of the
run, $N_c\le A_c$ for all $c$, and
\[
A_c=\ell M\,(1+nM)^{\,c-1}.
\]
In particular every exponent occurring has absolute value at most $B=A_n=\ell M(1+nM)^{n-1}$,
of bit-length $\log_2 B=O\!\bigl(n(b+\log n)+\log\ell\bigr)$.
\end{proposition}

\begin{proof}
\emph{Solving the recursion.} For $c\ge2$, $A_c-A_{c-1}=nM\,A_{c-1}$, so $A_c=(1+nM)A_{c-1}$
with $A_1=\ell M$, giving $A_c=\ell M(1+nM)^{c-1}$. Taking $\log_2$,
\[
\log_2 B=\log_2\ell+b+(n-1)\log_2(1+n2^{b})=O\!\bigl(\log\ell+n(b+\log n)\bigr).
\]

\emph{The bound $N_c\le A_c$.} For each level $c$, define $T_c$ = the total level-$c$ mass
ever \emph{created} during the run, counting the initial input as creation at time $0$ and
counting each newly inserted syllable's $|e|$ once. Clearly $N_c\le T_c$ at every moment
(present mass never exceeds created mass, since masses are only created, transferred, or
destroyed). We bound $T_c$.

By Lemma~\ref{lem:level}, level-$c$ mass is created only:
\begin{itemize}
\item[(i)] initially, by the input: at most $\ell$ syllables of exponent $<M$, contributing
$\le \ell M$;
\item[(ii)] by an iteration acting on a level-$j$ syllable with $j<c$ (the only iterations
that can create level-$c$ mass, again by Lemma~\ref{lem:level}, since a level-$c$ or higher
iteration never creates level-$<c$, hence to create level $c$ the acting level must be $<c$;
moreover same-level ($j=c$) iterations only create level-$c$ residuals that conserve $N_c$, as
shown below, so they do not increase $T_c$ on net — we account for net creation).
\end{itemize}

Let us quantify (ii). Consider a single iteration acting on a level-$j$ syllable $(j,e)$ with
$j<c$. It is either a swap ($c'>j$ on the prefix side) or a power reduction:
\begin{itemize}
\item Power reduction at level $j$ inserts $|q|\le|e|$ copies of $w_{j,j}^{\pm1}$; each copy
contributes to $N_c$ at most the absolute value of the level-$c$ exponent of $w_{j,j}$, which
is $<M$, and there is at most one level-$c$ syllable per copy. So this iteration adds at most
$|e|\cdot M$ to $T_c$.
\item Swap at level $j<k$ inserts $\mathrm{Conj}(g_k^{f},g_j^{\sigma})$ consisting of $|f|$
copies of $w_{j,k}^{\pm1}$; each copy adds at most $M$ to $T_c$ (one level-$c$ syllable of
exponent $<M$). So this iteration adds at most $|f|\cdot M$ to $T_c$, where $f$ is the exponent
of the prefix syllable being consumed, $|f|\le N_k\le T_k$.
\end{itemize}

Now sum over the whole run. We introduce one accounting unit per syllable that is ever
\emph{removed} from the prefix $P$ or \emph{reduced} as a power, and bound the level-$c$ mass
each such unit can create. Concretely, every iteration that creates level-$c$ mass is, by
Lemma~\ref{lem:level}, either a swap acting at some level $j<c$ or a power reduction at some
level $j<c$; in each case it removes from circulation either the prefix syllable $g_k^{f}$
(a swap, consuming $|f|$ units of level-$k$ mass, $k>j$) or the syllable $g_j^{e}$ (a power
reduction, consuming the part $|e|-e_0\ge r_j$ of its level-$j$ mass). By the two bullet
points, each such removed unit of mass generates at most $M$ units of mass at any one fixed
higher level, hence at most $M$ units of level-$c$ mass.

It remains to bound the total mass removed at each level. Consider level $j$. Mass at level
$j$ is created only as in (i)--(ii), so the total level-$j$ mass ever in circulation is exactly
$T_j$. Each unit of it is removed at most once (once a syllable's mass is consumed by a swap or
reduction it is gone, replaced only by mass at strictly higher levels). Therefore the total
mass removed at level $j$ over the whole run is at most $T_j$. Summing the per-removed-unit
contribution $\le M$ to level $c$, over all levels $j<c$ from which level-$c$ mass can be fed,
and recalling that a single removed unit may, through a chain of swaps, feed at most $n$
distinct higher target levels (one per level above $j$), we obtain the uniform bound
\[
T_c\;\le\;\ell M\;+\;nM\sum_{j<c}T_j ,
\]
valid simultaneously for every $c$. This is exactly the defining recursion of $A_c$. By
induction on $c$ (the right-hand side involves only $T_j$ with $j<c$), $T_c\le A_c$, and
therefore $N_c\le T_c\le A_c=\ell M(1+nM)^{c-1}\le B$ for all $c$. Every individual exponent
has absolute value at most the mass of its level, hence at most $B$.
\end{proof}

\section{Number of operations, termination, and total running time}\label{sec:steps}

\begin{proposition}\label{prop:steps}
Algorithm~\ref{alg:collect} performs at most
\begin{equation}\label{eq:totalsteps}
T_{\mathrm{op}}\;\le\;\ell\,(1+nMB)^{\,n}\;=\;\ell\cdot 2^{\,O(n\log B)}\;=\;\ell\cdot 2^{\,O(n^2(b+\log(n\ell)))}
\end{equation}
elementary operations \textup{(}placements, merges, power reductions, swaps\textup{)}; in
particular it terminates. Each operation costs $O\!\bigl(n\,\mathsf M(\log B)\bigr)$ bit
operations, where $\mathsf M(t)=\widetilde O(t)$ is the cost of multiplying $t$-bit integers
and $\log B=O(n(b+\log(n\ell)))$. Hence the total running time is
\[
\ell\cdot 2^{\,O(n^2(b+\log(n\ell)))}=(nM\ell)^{O(n^2)},
\]
proving parts \textup{(1)}--\textup{(2)} of Theorem~\ref{thm:main}.
\end{proposition}

\begin{proof}
\emph{Counting operations by level.} Let $\mathcal O_c$ be the number of iterations whose
acting syllable has level $c$. We claim
\begin{equation}\label{eq:Oc}
\mathcal O_c\;\le\; 2\,T_c\;\le\;2A_c\;\le\;2B,
\end{equation}
where $T_c$ is, as in Proposition~\ref{prop:exp}, the total level-$c$ mass ever created.
Indeed, each level-$c$ iteration permanently removes work at level $c$: a placement moves a
level-$c$ syllable into $P$, where it stays for the rest of the run (placed syllables are only
ever removed again by a \emph{higher}-or-equal-level swap that consumes them as the prefix tail
— but such a swap is counted at \emph{its} acting level, not $c$); a power reduction at level
$c$ removes $\ge2$ units of level-$c$ mass; a merge removes one of two level-$c$ syllables; a
swap with acting level $c$ transfers exactly one unit of level-$c$ mass leftward toward
placement. Thus, charging each level-$c$ iteration to either a unit of level-$c$ mass it moves
toward $P$ or a unit it destroys, and since the total level-$c$ mass ever created is $T_c$,
the number of level-$c$ iterations is at most twice $T_c$ (the factor $2$ covers a swap moving
a unit left followed later by the placement of that unit). This gives \eqref{eq:Oc}.

\emph{Summing over levels.} By \eqref{eq:Oc} and Proposition~\ref{prop:exp},
\[
T_{\mathrm{op}}=\sum_{c=1}^{n}\mathcal O_c\;\le\;\sum_{c=1}^{n}2A_c
=2\sum_{c=1}^{n}\ell M(1+nM)^{c-1}
=2\ell M\,\frac{(1+nM)^{n}-1}{nM}
\;\le\;2\ell\,(1+nM)^{n}.
\]
Since $(1+nM)^n\le(1+nMB)^n$, this yields \eqref{eq:totalsteps}. Because $B$, hence each $A_c$,
is finite (Proposition~\ref{prop:exp}), $T_{\mathrm{op}}$ is finite: the algorithm performs
finitely many operations and therefore halts.

To see that $T_{\mathrm{op}}<\infty$ truly forces termination (and not an infinite run with
finitely many ``operations''), note every iteration of the while-loop either (a) discards a
zero syllable, (b) is one of the four elementary operations counted above, in which case it is
charged in $T_{\mathrm{op}}$, or (c) is a placement. Placements append to $P$ and are counted
in the level sum; discards strictly shorten $U$ and create nothing. As $T_{\mathrm{op}}$ bounds
all non-discard iterations and discards cannot occur infinitely often between two non-discard
iterations (each discard removes a syllable from $U$ that was inserted by an earlier counted
operation, and only finitely many syllables are ever inserted, namely at most
$\sum_c \mathcal O_c\cdot n\le nT_{\mathrm{op}}$ of them), the total number of while-loop
iterations is at most $(n+1)T_{\mathrm{op}}<\infty$. Hence the loop terminates.

\emph{Cost per operation.} An elementary operation manipulates the relator word being inserted
($\le n$ syllables) and a constant number of other syllables, performing $O(n)$ integer
additions/subtractions and at most one division-with-remainder (in a power reduction), all on
integers of absolute value $\le B$ by Proposition~\ref{prop:exp}, i.e.\ of bit-length $\le
\log_2 B+1=O(n(b+\log(n\ell)))$. Each such arithmetic step costs $O(\mathsf M(\log B))=
\widetilde O(\log B)$ bit operations with fast integer arithmetic (or $O((\log B)^2)$ with
schoolbook arithmetic). So the per-operation cost is $O(n\,\mathsf M(\log B))$.

\emph{Total.} Multiplying the operation count $(n+1)T_{\mathrm{op}}$ by the per-operation cost,
\[
(n+1)\,T_{\mathrm{op}}\cdot O\!\bigl(n\,\mathsf M(\log B)\bigr)
=\ell\cdot 2^{O(n\log B)}\cdot \widetilde O(n^2\log B)
=\ell\cdot 2^{O(n^2(b+\log(n\ell)))},
\]
using $\log B=O(n(b+\log(n\ell)))$ and absorbing the polynomial factor into the exponential.
This is \eqref{eq:upper}.
\end{proof}

\begin{remark}
For a \emph{fixed} presentation ($n=O(1)$, $b=O(1)$, hence $M=O(1)$), \eqref{eq:upper} reads
$\ell^{O(1)}$: collection is polynomial in the input length for each fixed polycyclic
presentation, recovering the classical fact. The blow-up is entirely in the dependence on $n$.
\end{remark}

\section{Lower bound: the number of collection operations is exponential in $n$}\label{sec:lower}

The upper bound of Proposition~\ref{prop:steps} is exponential in $n$. We now show this is
genuine: there is a family of consistent polycyclic presentations with $b=1$ and inputs of
length $\mathrm{poly}(n)$ on which Algorithm~\ref{alg:collect} is forced to perform $2^{\Omega(n)}$
elementary operations. We are careful to claim only what is true. The exponential cost lies in
the \emph{operation count} of collection-from-the-left, which processes relator copies one at a
time; it does \emph{not} come from the size of the output, which on these instances is of
polynomial bit-length $O(n^2)$.

\subsection*{The construction}
Fix $n\ge2$ and let $J$ be the $n\times n$ lower-unitriangular Jordan block,
$J=I+S$ where $S$ is the nilpotent down-shift $S\,e_i=e_{i+1}$ ($e_{n+1}:=0$); equivalently $J$
has $1$'s on the diagonal and on the first \emph{sub}diagonal and $0$'s elsewhere, so
$J\,e_i=e_i+e_{i+1}$. Let
\[
H_n \;=\; \mathbb Z^{\,n}\rtimes_{J}\mathbb Z ,
\]
where the generator $t$ of $\mathbb Z$ acts on the free abelian group
$A=\langle a_1,\dots,a_n\rangle\cong\mathbb Z^{n}$ by $t^{-1}\mathbf v\,t=J\mathbf v$ (column
convention, $a_i\leftrightarrow e_i$), i.e.
\begin{equation}\label{eq:jordanrel}
a_i^{\,t}\;=\;a_i\,a_{i+1}\quad(1\le i\le n-1),\qquad a_n^{\,t}=a_n .
\end{equation}
This is a polycyclic group of Hirsch length $n+1$ (it is a finitely generated torsion-free
nilpotent group, being the standard unipotent extension; abelian-by-cyclic with both pieces
free abelian). 

\begin{lemma}\label{lem:Hpres}
Order the generators as $g_0=t$ \textup{(}level $0$\textup{)} and $g_i=a_i$ \textup{(}level
$i$, $1\le i\le n$\textup{)}. With this ordering, $H_n$ has a consistent polycyclic presentation
in which every relative order is $\infty$ \textup{(}no power relations\textup{)} and every
conjugation relator is a normal word in higher-level generators with all exponents in
$\{-1,0,1\}$; thus $b=1$ and $M=2$.
\end{lemma}

\begin{proof}
The $a_i$ commute pairwise, so for $1\le i<j$ the conjugation relators $a_j^{\,a_i}=a_j$,
$a_j^{\,a_i^{-1}}=a_j$ are trivial normal words. For the action of $t=g_0$ on $a_i=g_i$ the
relator $a_i^{\,t}=a_i a_{i+1}$ from \eqref{eq:jordanrel} is the normal word $g_i^{1}g_{i+1}^{1}$
(or $g_i^1$ if $i=n$), with exponents in $\{0,1\}$. The inverse action is
$a_i^{\,t^{-1}}$, the $i$-th column of $J^{-1}$. Since $J^{-1}=(I+S)^{-1}=\sum_{m\ge0}(-1)^mS^m$
and $S^m e_i=e_{i+m}$ (with $e_k=0$ for $k>n$), we get $J^{-1}e_i=\sum_{m\ge0}(-1)^m e_{i+m}$;
in our column convention this reads
\[
a_i^{\,t^{-1}}\;=\;a_i\,a_{i+1}^{-1}\,a_{i+2}\,a_{i+3}^{-1}\cdots a_n^{\pm1},
\]
again a normal word in $g_i,g_{i+1},\dots,g_n$ with exponents in $\{-1,+1\}$. Thus all relators
are normal words with $\{-1,0,1\}$-exponents, so $b=1$, $M=2$. Finally the presentation is
consistent: $H_n=A\rtimes\langle t\rangle$ realises it faithfully, and each element has a unique
expression $t^{a_0}a_1^{a_1}\cdots a_n^{a_n}$ ($a_0\in\mathbb Z$, $a_i\in\mathbb Z$), because
$A\cong\mathbb Z^n$ is free abelian with basis $a_1,\dots,a_n$ and the extension by
$\langle t\rangle\cong\mathbb Z$ splits; the normal words are exactly these expressions, in
bijection with $H_n$, which is consistency.
\end{proof}

\subsection*{The input and its normal form}
Take the input word
\begin{equation}\label{eq:input}
W_n \;=\; t^{-n}\,a_1\,t^{\,n},\qquad\text{of syllable length }\ell=2n+1.
\end{equation}

\begin{lemma}\label{lem:binom}
In $H_n$ the normal form of $W_n$ is
\[
W_n \;=\; a_1^{\binom n0}\,a_2^{\binom n1}\,a_3^{\binom n2}\cdots a_{n}^{\binom{n}{\,n-1}}
\;=\;\prod_{i=1}^{n} a_i^{\binom{n}{\,i-1}} .
\]
In particular it contains the exponent $\binom{n}{\lfloor n/2\rfloor}$, which satisfies
$\binom{n}{\lfloor n/2\rfloor}\ge 2^{n}/(n+1)$ and hence has bit-length $n-O(\log n)=\Theta(n)$.
The total exponent mass of the normal form is $\sum_{i=1}^n\binom{n}{i-1}=2^{n}$.
\end{lemma}

\begin{proof}
By definition of the action, conjugation by $t$ acts on the column vector of $a$-exponents by
the matrix $J$: writing an element $\prod_i a_i^{v_i}$ of $A$ as the column vector
$\mathbf v=(v_i)$, we have $t^{-1}\bigl(\prod_i a_i^{v_i}\bigr)t=\prod_i a_i^{(J\mathbf v)_i}$,
because the relation \eqref{eq:jordanrel} is exactly $t^{-1}a_i t=a_ia_{i+1}$, i.e.\ $J e_i$,
and the $a_i$ commute. Iterating $k$ times,
\[
t^{-k}\,a_1\,t^{k}\;=\;\prod_{i=1}^{n} a_i^{(J^{k}e_1)_i}\;=\;\prod_{i=1}^{n} a_i^{(J^k)_{i,1}} .
\]
We claim $(J^k)_{i,1}=\binom{k}{\,i-1}$ for all $k\ge0$ and $1\le i\le n$. This is the standard
formula for powers of a unipotent Jordan block: since $J=I+S$ with $S$ nilpotent and
$S^{m}e_1=e_{1+m}$, the binomial theorem gives $J^{k}e_1=(I+S)^{k}e_1=\sum_{m=0}^{k}\binom km
S^{m}e_1=\sum_{m=0}^{k}\binom km e_{1+m}$, whose $i$-th coordinate is $\binom{k}{\,i-1}$
(the term $m=i-1$), and $0$ once $i-1>k$. Taking $k=n$ gives
$t^{-n}a_1 t^{n}=\prod_{i=1}^n a_i^{\binom{n}{i-1}}$, which is already a normal word (increasing
levels $1<2<\cdots<n$, all relative orders infinite), hence is the normal form of $W_n$. The
central binomial coefficient bound $\binom{n}{\lfloor n/2\rfloor}\ge 2^n/(n+1)$ is the standard
estimate $\binom{n}{\lfloor n/2\rfloor}=\max_i\binom ni\ge\frac1{n+1}\sum_i\binom ni=2^n/(n+1)$.
The total mass is $\sum_{i=0}^{n-1}\binom ni=\sum_{i=0}^n\binom ni-\binom nn=2^n-1$ over levels
$1,\dots,n$; including the contribution we used we record $\sum_{i=1}^n\binom n{i-1}=2^n$.
\end{proof}

\noindent
\emph{Verification.} We checked Lemma~\ref{lem:binom} symbolically for $2\le n\le 6$ by running
Algorithm~\ref{alg:collect} on $W_n$ in $H_n$: the output is exactly
$a_1 a_2^{\binom n1}\cdots a_n^{\binom n{n-1}}$, e.g.\ for $n=5$ it is
$a_1\,a_2^{5}a_3^{10}a_4^{10}a_5^{5}$.

\subsection*{The operation lower bound}
The key point is a conservation argument: every unit of exponent mass at levels $\ge2$ in the
\emph{output} had to be \emph{created} by Algorithm~\ref{alg:collect}, and each unit of created
mass costs at least one elementary operation.

\begin{lemma}[Created syllables lower-bound operations]\label{lem:opsmass}
Run Algorithm~\ref{alg:collect} on any input, and assume \textup{(}as for $H_n$\textup{)} that
\emph{every relator exponent lies in $\{-1,0,1\}$}, so that every syllable created during the
run carries an exponent of absolute value exactly $1$. For a level $c\ge1$ let $S_c^{\mathrm{tot}}$
be the total number of level-$c$ syllables that ever appear during the run, counting input
syllables and each newly inserted syllable once. Then the number of elementary operations is at
least $\sum_{c\ge1} S_c^{\mathrm{tot}}-(\text{number of input syllables})$, and in particular at
least the number of level-$c$ syllables created \emph{after} time $0$, summed over $c\ge1$.
Consequently, writing $I_c$ for the input level-$c$ exponent mass and $N_c^{\mathrm{out}}$ for
the output level-$c$ exponent mass, the number of operations is at least
\[
\sum_{c\ge1}\max\!\bigl(0,\,N_c^{\mathrm{out}}-I_c\bigr).
\]
\end{lemma}

\begin{proof}
Every iteration of the while-loop pops exactly one syllable from $U$; thus the number of
elementary operations equals the number of syllables ever popped, which is at least the number
of syllables ever \emph{created} minus those never popped. Every created syllable except possibly
the final placed ones is popped (placed syllables sit in $P$, but a placement is itself an
operation that pops the syllable from $U$ first); hence the number of operations is at least the
total number of syllables ever created after time $0$, i.e.\ $\sum_c S_c^{\mathrm{tot}}-(\text{\#
input syllables})$.

By Lemma~\ref{lem:level}, level-$c$ syllables are created only in the input or, one per syllable
of an inserted relator copy, during a swap or power reduction acting at some level $j<c$ (the
only events producing level-$c$ syllables). Under the hypothesis that all relator exponents are
$\pm1$ or $0$, every such created syllable has exponent of absolute value $1$. Therefore each
created level-$c$ syllable contributes exactly one unit of level-$c$ mass at the instant of its
creation. The output level-$c$ mass $N_c^{\mathrm{out}}$ is built up from these unit
contributions (and from the at most $I_c$ units present in the input) by additive merging only;
mass is never increased except by creating a $\pm1$ syllable. Hence at least
$N_c^{\mathrm{out}}-I_c$ level-$c$ syllables of exponent $\pm1$ must have been created after time
$0$ to account for the output mass at level $c$. Summing over $c$ and using that each created
syllable is charged to a distinct pop, the number of operations is at least
$\sum_{c\ge1}\bigl(N_c^{\mathrm{out}}-I_c\bigr)$ whenever every summand is nonnegative.
\end{proof}

\begin{proposition}\label{prop:lower}
For the family $(H_n,W_n)$ of Lemma~\ref{lem:Hpres}--\ref{lem:binom}, with $b=1$, $M=2$ and
input length $\ell=2n+1=\mathrm{poly}(n)$, Algorithm~\ref{alg:collect} performs at least
$2^{n}-n$ elementary operations. Hence the operation count of collection-from-the-left is
$2^{\Omega(n)}$, exponential in the Hirsch length, even when $b=O(1)$ and $\ell=\mathrm{poly}(n)$.
\end{proposition}

\begin{proof}
The input $W_n=t^{-n}a_1 t^n$ has all of its $a$-mass at level $1$ (one syllable $a_1$) and
none at levels $\ge2$: thus $I_1=1$ and $I_c=0$ for $c\ge2$. By Lemma~\ref{lem:binom} the output
mass at level $c$ is $N_c^{\mathrm{out}}=\binom{n}{c-1}$. By Lemma~\ref{lem:opsmass} the number
of operations is at least
\[
\sum_{c\ge1}\max\!\bigl(0,\,N_c^{\mathrm{out}}-I_c\bigr)
=\bigl(\binom n0-1\bigr)+\sum_{c=2}^{n}\binom{n}{c-1}
=0+\sum_{i=1}^{n-1}\binom ni
=2^{n}-2 .
\]
Thus Algorithm~\ref{alg:collect} performs at least $2^n-2\ge 2^n-n$ elementary operations on
$(H_n,W_n)$. Since $\ell=2n+1$ and $b=1$, this is $2^{\Omega(n)}$ with $\ell=\mathrm{poly}(n)$
and $b=O(1)$, establishing the necessity of exponential dependence on $n$ in the operation
count.
\end{proof}

\noindent\emph{Verification.} Running Algorithm~\ref{alg:collect} on $(H_n,W_n)$ and counting
iterations gives $14,33,78,197,546,1663,5480$ operations for $n=2,\dots,8$, all $\ge 2^n-2$
($2,6,14,30,62,126,254$) as the proposition guarantees, confirming the bound is correct (and
in fact the true count grows somewhat faster than $2^n$).

\begin{remark}[What the lower bound does and does not say]
Proposition~\ref{prop:lower} bounds the \emph{work} of collection-from-the-left, not the output
size. On $(H_n,W_n)$ the normal form has $n$ syllables whose exponents are binomial coefficients
$\binom{n}{i}\le 2^n$, i.e.\ of bit-length $\le n$; the whole output therefore occupies
$O(n^2)$ bits, which is \emph{polynomial} in $n$. The exponential blow-up is intrinsic to the
algorithm: it expands the conjugation/power relators one copy at a time
(lines~\ref{ln:swapprepend}, \ref{ln:powerprepend}), so moving a prefix syllable $g_k^f$ across
costs $\Theta(|f|)$ operations, and the accumulated mass $2^n$ that must be carried forces
$2^{\Omega(n)}$ operations. This matches, in its $n$-dependence, the upper bound
$T_{\mathrm{op}}=\ell\cdot 2^{O(n^2(b+\log(n\ell)))}$ of Proposition~\ref{prop:steps}
\textup{(}both are $2^{\Theta(\mathrm{poly}(n))}$ with $b,\ell=O(1)$\textup{)}, and it shows the
classical fact that collection is polynomial \emph{for each fixed presentation} cannot be made
uniform in the Hirsch length.
\end{remark}

\begin{remark}[On the role of exponent magnitudes]
The exponent-size bound \eqref{eq:expbound} shows that every integer manipulated by the
algorithm has bit-length $O(n(b+\log(n\ell)))$, i.e.\ \emph{polynomial} in $n,b,\log\ell$.
Thus the exponential cost is genuinely combinatorial — an exponential \emph{number} of
operations on polynomially-sized integers — rather than arithmetic on astronomically large
numbers. This is why a careful collection implementation that replaces unary relator-copy
insertion by repeated squaring of the conjugating action can, on nilpotent groups, reduce the
$n$-dependence; the present analysis pertains to collection-from-the-left exactly as specified
in Algorithm~\ref{alg:collect}, the classical algorithm whose worst-case the problem asks for.
\end{remark}

\section{Conclusion}

Combining Propositions~\ref{prop:exp}, \ref{prop:steps} and \ref{prop:lower} proves
Theorem~\ref{thm:main}. The worst-case complexity of collection-from-the-left for computing
the normal form of a length-$\ell$ word in a polycyclic group of Hirsch length $n$, presented
by a consistent polycyclic presentation with exponent bit-bound $b$, is
\[
(nM\ell)^{O(n^2)}\;=\;\ell\cdot 2^{\,O(n^2(b+\log(n\ell)))},\qquad M=2^b,
\]
with all intermediate and final exponents of bit-length $O\!\bigl(n(b+\log(n\ell))\bigr)$. This
is polynomial in the input length $\ell$ and in the exponent magnitude $M=2^b$ for every fixed
Hirsch length $n$ (recovering the classical fact that collection is polynomial for a fixed
presentation), while the number of elementary collection operations — and hence the running
time — is exponential in the Hirsch length $n$; the latter dependence is necessary, as
witnessed by the unipotent groups $H_n=\mathbb Z^n\rtimes_J\mathbb Z$ of
Proposition~\ref{prop:lower}, on which collection-from-the-left performs at least $2^n-n$
elementary operations although $b=1$, $\ell=\mathrm{poly}(n)$, and the output has only
polynomial bit-size.
\end{document}
